\section{Introduction}
Automating routine computer tasks with autonomous agents is a long standing goal of artificial intelligence research~\citep{franklin1996agent,jennings1998roadmap}. To achieve this, we need agents that can navigate computers effectively, process visual and textual inputs, handle high-level natural language instructions, and execute actions to achieve desired goals. As digital interfaces today are primarily built for human eyes, effective visual understanding is necessary for many routine computer tasks. For example, humans frequently perform tasks on the web which involve visual references, such as ``Help me order a green polo shirt from Amazon,'' or rely on pictures rather than text to communicate. %Productive work done on the computer is also often explicitly visual, e.g., working with Microsoft Excel sheets, creating presentations in Google Slides, or performing creative work in Adobe Photoshop. 
However, many agent benchmarks today focus on text-based tasks, neglecting the evaluation (and consequently the development) of multimodal agents.
\begin{figure*}[t]
\centering
\includegraphics[width=1.0\linewidth]{figures/overview.pdf}
% \vspace{-0.3in}
\caption{\BenchmarkName{} is a benchmark suite of 910 realistic, visually grounded tasks on self-hosted web environments that involve web navigation and visual understanding.
}
% \vspace{-0.1in}
\label{fig:task_overview}
\end{figure*}
To address this gap, we propose \BenchmarkName{}~(Fig.~\ref{fig:task_overview}), a benchmark suite designed to rigorously assess and advance the visual and textual capabilities of autonomous agents. 
\BenchmarkName{} builds off the WebArena~\citep{zhou2023webarena} framework, leveraging reproducible self-hosted environments and execution-based evaluations. 
%\BenchmarkName{} introduces a set of unique visual tasks that replicate those performed by humans on realistic real-world websites. 
\BenchmarkName{} introduces a set of unique tasks that emphasize integrating visual understanding with language processing, closely simulating human interaction with modern computing interfaces. 
% By introducing \BenchmarkName{}, we aim to bridge the existing gap in the evaluation of multimodal autonomous agents and to push the boundaries of AI-driven task automation. %We envision a future where AI agents can not only process textual information but also effectively navigate and interact with visual settings (e.g., digital screens), truly augmenting human productivity and creativity. 
Our contributions are summarized as follows:

\begin{itemize}
    \item We introduce \BenchmarkName{}, a set of 910 realistic tasks over three diverse web environments: Classifieds, Shopping, and Reddit. The Classifieds environment is a new contribution with real world data, while the Shopping and Reddit environments are inherited from WebArena. All tasks we introduce are \textit{visually grounded}, and require visual understanding of webpage content to effectively solve (while WebArena does not). 25.2\% of the tasks also include images as input (Fig.~\ref{fig:task_overview}), and require understanding interleaved image-text inputs. %\BenchmarkName{} will be publicly released under an Apache 2.0 license. % to correctly process and attempt. Our code, tasks, and baselines are publicly released.\footnote{\url{url_removed_for_review}}
    % \item To enable benchmarking of realistic open-ended tasks, we introduce new automated visual evaluation metrics: (1) VQA evaluation, where we query a vision-language model to assess whether the completed task satisfies certain visual criteria (e.g., whether the page contains an image of sushi) and (2) image similarity match, which checks whether a page contains the exact same image as a given groundtruth image.
    \item We extensively benchmark the autonomous capabilities of state-of-the-art (SOTA) large language models (LLM) and vision-language models (VLMs), demonstrating that strong VLMs outperform text-based LLMs. The best VLM agents achieve a success rate of 16.4\% on \BenchmarkName{}, which is still significantly below human performance of 88.7\%. %Our results also highlight a large gap between API-based and open sourced VLM agents.
    \item We propose a new VLM agent inspired by Set-of-Marks prompting~\cite{yang2023som}, simplifying the action space of the model. We show that this model substantially outperforms other baseline LLM agents, especially on sites that are more visually complex.
\end{itemize}

